% problem_id: S047-aif3604-P2
% source_id: S047 (Corrigendum to “A simpler proof of toroidalization of morphisms from 3-folds to surfaces” (Steven Dale Cutkosky; Ann. Inst. Fourier, Grenoble 74, 6 (2024) 2505–2521; DOI 10.5802/aif.3604; 18 PDF pages). The corrigendum repairs Ann. Inst. Fourier 63 (2013) 865–922 — the paper [3] of its own bibliography.)
% source_file: aif3604.pdf
% statement_locator: printed p. 4
% proof_locator: printed p. [5, 6, 7, 8]
% representation: PDF; transcribed from rendered page images (never from the text layer)
% collected_by: carrier rule (the headline theorem was an assembly; this is the step it relies on)
% review_status: PENDING independent review (single delegated judgement)
% status: complete (statement + author proof verbatim + reason + locators)
%
%% problem_id: P2
%% collected unit: Lemma 3.6 (as REVISED by the errata) together with its complete proof
%% paper: Steven Dale Cutkosky, "Corrigendum to `A simpler proof of toroidalization of morphisms
%%        from 3-folds to surfaces'", Ann. Inst. Fourier, Grenoble 74, 6 (2024) 2505-2521,
%%        DOI 10.5802/aif.3604 (18 PDF pages). The errata corrects Ann. Inst. Fourier 63 (2013)
%%        865-922, whose Lemma 3.6 misses a case; this errata exists precisely to supply that case.
%% source_locator: /disks/sata1/yupeng/human-proof-corpus/source-cache/registry-sources/numdam/aif3604.pdf
%%
%% PAGE MAP (the printed page number is 2503 + the PDF page number; PDF page 1 is the journal cover,
%% which carries no printed number, and printed page 2505 = PDF page 2):
%%   context for the counted item (tail of the revised Definition 3.3, items (3),(4)/(A)): PDF p. 4
%%   counted STATEMENT of Lemma 3.6:  PDF pp. 4-5  = printed pp. 2507-2508
%%   counted PROOF of Lemma 3.6:      PDF pp. 5-8  = printed pp. 2508-2511 (ends with the end-of-proof box)
%%   AUXILIARY block (NOT part of the counted unit): the inserted detailed proof of the case-(B)
%%     statement of the revised Lemma 3.7, PDF pp. 10-12 = printed pp. 2513-2515; transcribed only
%%     because it is the continuation of the very same new case through the next errata item.
%%
%% Reading method: the PDF was rendered on the host with
%%   pdftoppm -f <p> -l <p> -r 150 -png aif3604.pdf <pfx>
%% and every formula below was additionally re-read from 300 dpi half-page renders and from targeted
%% 600 dpi crops (pdftoppm -r 600 -x -y -W -H). The pdftotext text layer of this file is usable only
%% for locating: it drops the bars and hats (it prints $\bar a_i$ as "ai", $\bar v$ as "v",
%% $\tilde y$ as "ye", $\hat x_1$ as "x1", $J_{\mathfrak p}$ as "Jp"), so it was never used as source
%% text. Every line transcribed below was READ FROM THE PAGE IMAGES.
%%
%% Transcription policy: verbatim. Nothing rewritten, simplified, completed or silently corrected.
%% Displayed formulas are reproduced exactly as displayed, including their printed slips; each slip is
%% flagged with a %% [NOTE: ...] comment and is NOT repaired in the body. Line breaks of the original
%% prose are not reproduced. No passage was illegible, so no [ILLEGIBLE: ...] marker occurs.

\documentclass[11pt]{article}
\usepackage{amsmath}
\usepackage{amssymb}
\usepackage{amsthm}

\theoremstyle{plain}
\newtheorem{lemma}{Lemma}

%% \Bbbk (blackboard-bold k) is already provided by amssymb; the paper's field symbol is used as printed.
\newcommand{\Oh}[1]{\widehat{\mathcal{O}}_{#1}}
\newcommand{\Spec}{\operatorname{Spec}}

\begin{document}

%% =====================================================================
%% CONTEXT (not counted): revised Definition 3.3, items (3) and (4)/(A), and the
%% instruction line that introduces the counted item.  PDF p. 4 = printed p. 2507.
%% Transcribed because the counted lemma is stated for morphisms that are "2-prepared"
%% with $m=\sigma_D(p)+1$, and because item (A) is the new case created by the errata.
%% =====================================================================

%% --- page 4 (printed 2507) ---

(3) $p$ is a 1-point, and we have an expression (2.1) with
\begin{equation}
F = \tau_0 z^m + \tau_2 x^{r_2} z^{m-2} + \cdots + \tau_{m-1} x^{r_{m-1}} z + x^t \Omega
\tag{3.7}
\end{equation}
where $\tau_0 \in \Oh{X,p}$ is a unit, $\tau_i \in \Oh{X,p}$ are units (or zero) for $2 \leqslant i \leqslant m-1$, $\Omega \in \Oh{X,p}$, $\tau_i \neq 0$ for some $i$ with $2 \leqslant i \leqslant m-1$ and $t > \omega(m, r_2, \ldots, r_{m-1})$ (where we set $r_i = 0$ if $\tau_i = 0$). Further, $r_i > 0$ if $\tau_i \neq 0$.

(4) $p$ is a 1-point, and we have an expression (2.1) with
\begin{equation}
F = \tau_0 z^m + \tau_2 x^{r_2} y z^{m-2} + \cdots + \tau_{m-1} x^{r_{m-1}} y z + \tau_m x^{r_m} y + x^t \Omega
\tag{A}
\end{equation}
where $\tau_0 \in \Oh{X,p}$ is a unit, $\tau_i \in \Oh{X,p}$ are units (or zero) for $2 \leqslant i \leqslant m$, $\Omega \in \Oh{X,p}$, $\tau_i \neq 0$ for some $i$ with $2 \leqslant i \leqslant m$ and $t > \omega'(m, r_2, \ldots, r_m)$ (where we set $r_i = 0$ if $\tau_i = 0$). Further, $r_i > 0$ if $\tau_i \neq 0$.

$X$ is 3-prepared if $X$ is 3-prepared for all $p \in X$.

Page 881, line 5: ``Lemma'' 3.5.

Pages 880--882: Lemma 3.6 and its proof should be changed as follows.

%% =====================================================================
%% COUNTED STATEMENT.  PDF pp. 4-5 = printed pp. 2507-2508.
%% =====================================================================

\begin{lemma}
Suppose that $X$ is 2-prepared with respect to $f : X \to S$. Suppose that $p \in D$ is a 1-point with $m = \sigma_D(p) + 1 > 1$. Let $u, v$ be permissible parameters for $f(p)$ and $x, y, z$ be permissible parameters for $D$ at $p$ such that a form (3.1) holds at $p$. Let $U$ be an \'etale cover of an affine neighborhood of $p$ such that $x, y, z$ are uniformizing parameters on $U$. Let $C$ be the curve in $U$ which has local equations $x = y = 0$ at $p$.

Let $T_0 = \Spec(\Bbbk[x, y])$, $\Lambda_0 : U \to T_0$. Then there exists a sequence of quadratic transforms $T_1 \to T_0$ such that if $U_1 = U \times_{T_0} T_1$ and $\psi_1 : U_1 \to U$ is the induced sequence of blow ups of sections over $C$, $\Lambda_1 : U_1 \to T_1$ is the projection, then $U_1$ is 2-prepared with respect to $f \circ \psi_1$ at all $p_1 \in \psi_1^{-1}(p)$. Further, for every point $p_1 \in \psi_1^{-1}(p)$, there exist regular parameters $x_1, y_1$ in $\Oh{T_1,\Lambda_1(p_1)}$ such that $x_1, y_1, z$ are permissible parameters at $p_1$, and there exist regular parameters $\tilde{x}_1, \tilde{y}_1$ in $\mathcal{O}_{T_1,\Lambda_1(p_1)}$ such that if $p_1$ is a 1-point, $x_1 = \alpha(\tilde{x}_1, \tilde{y}_1)\tilde{x}_1$ where $\alpha(\tilde{x}_1, \tilde{y}_1) \in \Oh{T_1,\Lambda_1(p_1)}$ is a unit series and $y_1 = \beta(\tilde{x}_1, \tilde{y}_1)$ with $\beta(\tilde{x}_1, \tilde{y}_1) \in \Oh{T_1,\Lambda_1(p_1)}$, and if $p_1$ is a 2-point, then $x_1 = \alpha(\tilde{x}_1, \tilde{y}_1)\tilde{x}_1$ and $y_1 = \beta(\tilde{x}_1, \tilde{y}_1)\tilde{y}_1$, where $\alpha(\tilde{x}_1, \tilde{y}_1), \beta(\tilde{x}_1, \tilde{y}_1) \in \Oh{T_1,\Lambda_1(p_1)}$ are unit series.
We have one of the following forms:
%% [NOTE: printed p. 2507 prints the sheaf symbol for $\tilde x_1,\tilde y_1$ WITHOUT the hat
%%  ("in $\mathcal{O}_{T_1,\Lambda_1(p_1)}$") while the line just above it prints it WITH the hat for
%%  $x_1,y_1$ ("in $\widehat{\mathcal{O}}_{T_1,\Lambda_1(p_1)}$"). Reproduced as printed.]
\end{lemma}

%% --- page 5 (printed 2508) ---

(1) $p_1$ is a 2-point, and we have an expression (2.2) with
\begin{equation}
F = \tau z^m + \bar{a}_2(x_1, y_1) x_1^{r_2} y_1^{s_2} z^{m-2} + \cdots + \bar{a}_{m-1}(x_1, y_1) x_1^{r_{m-1}} y_1^{s_{m-1}} z + \bar{a}_m x_1^{r_m} y_1^{s_m}
\tag{3.10}
\end{equation}
where $\tau \in \Oh{U_1,p_1}$ is a unit, $\bar{a}_i(x_1, y_1) \in \Bbbk[[x_1, y_1]]$ are units (or zero) for $2 \leqslant i \leqslant m-1$, $\bar{a}_m = 0$ or $1$ and if $\bar{a}_m = 0$, then $\bar{a}_{m-1} \neq 0$. Further, $r_i + s_i > 0$ whenever $\bar{a}_i \neq 0$ and $(r_m + c)b - (s_m + d)a \neq 0$.

(2) $p_1$ is a 1-point, and we have an expression (2.1) with
\begin{equation}
F = \tau z^m + \bar{a}_2(x_1, y_1) x_1^{r_2} z^{m-2} + \cdots + \bar{a}_{m-1}(x_1, y_1) x_1^{r_{m-1}} z + x_1^{r_m} y_1
\tag{3.11}
\end{equation}
where $\tau \in \Oh{U_1,p_1}$ is a unit, $\bar{a}_i(x_1, y_1) \in \Bbbk[[x_1, y_1]]$ are units (or zero) for $2 \leqslant i \leqslant m-1$. Further, $r_i > 0$ (whenever $\bar{a}_i \neq 0$).

(3) $p_1$ is a 1-point, and we have an expression (2.1) with
\begin{equation}
F = \tau z^m + \bar{a}_2(x_1, y_1) x_1^{r_2} z^{m-2} + \cdots + \bar{a}_{m-1}(x_1, y_1) x_1^{r_{m-1}} z + x_1^t y_1 \Omega
\tag{3.12}
\end{equation}
where $\tau \in \Oh{U_1,p_1}$ is a unit, $\bar{a}_i(x_1, y_1) \in \Bbbk[[x_1, y_1]]$ are units (or zero) for $2 \leqslant i \leqslant m-1$ and $r_i > 0$ whenever $\bar{a}_i \neq 0$. We also have $t > \omega(m, r_2, \ldots, r_{m-1})$. Further, $\bar{a}_i \neq 0$ for some $2 \leqslant i \leqslant m-1$ and $\Omega \in \Oh{U_1,p_1}$.

(4) $p_1$ is a 1-point, and we have an expression (2.1) with
\begin{equation}
F = \tau z^m + \bar{a}_2(x_1, y_1) x_1^{r_2} y_1 z^{m-2} + \cdots + \bar{a}_{m-1}(x_1, y_1) x_1^{r_{m-1}} y_1 z + x_1^{r_m} y_1 + x_1^t y_1^2 \Omega
\tag{B}
\end{equation}
where $\tau \in \Oh{U_1,p_1}$ is a unit, $\bar{a}_i(x_1, y_1) \in \Bbbk[[x_1, y_1]]$ are units (or zero) for $2 \leqslant i \leqslant m$ and $r_i > 0$ whenever $\bar{a}_i \neq 0$. We also have $t > \omega'(m, r_2, \ldots, r_m)$. Further, $\bar{a}_i \neq 0$ for some $2 \leqslant i \leqslant m$ and $\Omega \in \Oh{U_1,p_1}$.

%% =====================================================================
%% COUNTED PROOF.  PDF pp. 5-8 = printed pp. 2508-2511.
%% =====================================================================

\noindent\emph{Proof.} --- Let $\bar{p} = \Lambda_0(p)$. Let $T = \{i \mid a_i(x, y) \neq 0$ and $2 \leqslant i < m\}$. There exists a sequence of blow ups $\varphi_1 : T_1 \to T_0$ of points over $\bar{p}$ such that at all points $q \in \psi_1^{-1}(p)$, we have permissible parameters $x_1, y_1, z$ such that $x_1, y_1$ are regular parameters in $\Oh{T_1,\Lambda_1(q)}$ and we have that $ug$ is a monomial in $x_1$ and $y_1$ times a unit in $\Oh{T_1,\Lambda_1(q)}$, where $g = \prod_{i \in T} a_i(x, y)$.

Suppose that $a_m(x, y) \neq 0$. Let $\bar{v} = x^b a_m(x, y)$ if (2.1) holds and $\bar{v} = x^c y^d a_m(x, y)$ if (2.2) holds. We have $\bar{v} \notin \Bbbk[[x]]$ (respectively $\bar{v} \notin \Bbbk[[x^a y^b]]$). Then by Lemma 3.5 applied to $u, \bar{v}$, we have that there exists a further sequence of blow ups $\varphi_2 : T_2 \to T_1$ of points over $\bar{p}$ such that at all points $q \in (\psi_1 \circ \psi_2)^{-1}(p)$, we have permissible parameters $x_2, y_2, z$ such that $x_2, \bar{y}_2$

%% --- page 6 (printed 2509) ---

are regular parameters in $\Oh{T_2,\Lambda_2(q)}$ such that $ug = 0$ is a SNC divisor and either
\begin{equation}
u = x_2^{\bar{a}}, \bar{v} = \bar{P}(x_2) + x_2^{\bar{b}} \bar{y}_2^{\bar{c}}
\tag{c1}
\end{equation}
with $\bar{c} > 0$ or
\begin{equation}
u = (x_2^{\bar{a}} \bar{y}_2^{\bar{b}})^t, \bar{v} = \bar{P}(x_2^{\bar{a}} \bar{y}_2^{\bar{b}}) + x_2^{\bar{c}} \bar{y}_2^{\bar{d}}
\tag{c2}
\end{equation}
where $\bar{a}\bar{d} - \bar{b}\bar{c} \neq 0$.

At a 1-point or 2-point $p_2 \in U_2$ above $p$, we ask if there is an expression (2.1) or (2.2) such that $F$ has an expression
\begin{equation}
F = \tau z^m + \bar{a}_2(x_2, y_2) x_2^{r_2} y_2^{s_2} z^{m-2} + \cdots + \bar{a}_{m-1}(x_2, y_2) x_2^{r_{m-1}} y_2^{s_{m-1}} z + \bar{a}_m x_2^{r_m} y_2^{s_m}
\tag{3.13}
\end{equation}
where $\tau$ is a unit series, $\bar{a}_m = 0$ or $1$ and $\bar{a}_i$ are unit series (or zero) for $2 \leqslant i < m$.

At a 2-point $p_2 \in U_2$ above $p$ we always have an expression (2.2) such that (3.13) holds, since $ug = 0$ is a SNC divisor and by (c2).

Suppose that $p_2 \in U_2$ is a 1-point above $p$. Let $x_2, \bar{y}_2$ be the regular parameters of (c1) in $\Oh{T_2,\Lambda_2(p_2)}$. There exists $\tilde{y} \in \Oh{T_2,\Lambda_2(p_2)}$ such that $x_2, \tilde{y}$ are regular parameters and $ug = \bar{\alpha} x_2^{\bar{a}} \tilde{y}^{\bar{b}}$ where $\bar{\alpha}$ is a unit. Since $x_2, \bar{y}_2$ is a regular system of parameters in $\Oh{T_2,\Lambda_2(p_2)} \cong k[[x_2, \bar{y}_2]]$, we have that $\operatorname{ord} \tilde{y}(0, \bar{y}_2) = 1$. Thus by the Weierstrass Preparation Theorem, there exists a unit $\bar{\tau} \in k[[x_2, \bar{y}_2]]$ and nonunit series $\varphi(x_2) \in k[[x_2]]$ such that $\tilde{y} = \bar{\tau}(\bar{y}_2 + \varphi(x_2))$. Substituting into $F$, we obtain an expression (2.1) with
\begin{equation}
F = \tau z^m + \bar{a}_2(x_2, \bar{y}_2) x_2^{r_2} (\bar{y}_2 + \varphi(x_2))^{s_2} z^{m-2} + \cdots + \bar{a}_{m-1}(x_2, \bar{y}_2) x_2^{r_{m-1}} (\bar{y}_2 + \varphi(x_2))^{s_{m-1}} z + \bar{a}_m x_2^{r_m} \bar{y}_2^{s_m}
\tag{C}
\end{equation}
where $\tau$ is a unit series, $\bar{a}_m = 0$ or $1$, $s_m \geqslant 1$ and $\bar{a}_i$ are unit series (or zero) for $2 \leqslant i \leqslant m$, $\varphi \in k[[x_2]]$ is a series with $0 \leqslant r := \operatorname{ord} \varphi \leqslant \infty$, and $x_2, \bar{y}_2, z$ are permissible parameters with $\bar{y}_2 \in \Oh{T_2}$.

We will show that after a finite number of blow ups of points $T_3 \to T_2$, we have an expression (2.2) with (3.13) for all 2-points $p_3 \in U_3$ above $p$ and an expression (2.1) with (3.13) for all 1-points $p_3 \in U_3$ above p.

Suppose that $p_2 \in U_2$ is a 1-point. If $r = 0$ or $\infty$, $\bar{a}_m = 0$ or $s_m = 1$, then after a permissible change of variables we have a form (3.13). In fact, if a form (C) holds at $p_2$ with $s_m = 1$, then we set $\tilde{y}_2 = \bar{y}_2 + \varphi(x_2)$ and write
\[
v = P(x_2) + x_2^b F = P(x_2) + \bar{a}_m x_2^{b+r_m} \varphi(x_2) + x_2^b F'
\]

%% --- page 7 (printed 2510) ---

where
\[
F' = \tau z^m + \bar{a}_2 x_2^{r_2} \tilde{y}_2^{s_2} z^{m-2} + \cdots + \bar{a}_{m-1} x_2^{r_{m-1}} \tilde{y}_2^{s_{m-1}} z + \bar{a}_m x_2^{r_m} \tilde{y}_2
\]
%% [NOTE: in this display the last term is printed "$\bar a_m x_2^{r_m}\tilde y_2$", i.e. WITHOUT the
%%  exponent $s_m$ that form (3.13) carries on $y_2$; the preceding terms do carry their exponents.
%%  Reproduced as printed.]
to obtain the form (3.13).

Suppose that a form (3.13) does not hold at $p_2$. Then $0 < r = \operatorname{ord} \varphi < \infty$, $\bar{a}_m = 1$ and $s_m > 1$ in (C). Let $T_3 \to T_2$ be the blow up of $\Lambda_2(p_2) \in T_2$ with induced blow up $\psi_3 : U_3 \to U_2$. Suppose $p_3 \in \psi_3^{-1}(p_2)$. We have permissible parameters $x_3, y_3, z$ in $\Oh{U_3,p_3}$ such that one of the following cases holds:
\begin{align*}
\text{(a)}\quad & x_2 = x_3,\ \bar{y}_2 = x_3(y_3 + \alpha) \text{ with } 0 \neq \alpha \in \Bbbk. \\
\text{(b)}\quad & x_2 = x_3,\ \bar{y}_2 = x_3 y_3. \\
\text{(c)}\quad & x_2 = x_3 y_3,\ \bar{y}_2 = y_3.
\end{align*}

Suppose that (a) holds at $p_3$. We have
\[
\bar{y}_2^{s_m} = x_3^{s_m}(y_3 + \alpha)^{s_m} = x_3^{s_m}(\alpha^{s_m} + \mu(y_3)y_3)
\]
where $\mu(y_3)$ is a unit series.

If $\alpha + \frac{1}{x_3}\varphi(x_3)$ is a unit series, then setting $\tilde{y}_3 = \mu(y_3)y_3$, we have
\[
v = P(x_2) + x_2^b F = P(x_3) + \bar{a}_m \alpha^{s_m} x_3^{b+r_m+s_m} F'
\]
where
\[
F' = \tau z^m + a'_2 x_2^{r_2+s_2} + \cdots + a'_{m-1} x_3^{r_{m-1}+s_{m-1}} + \bar{a}_m x_3^{r_m+s_m} \tilde{y}_3,
\]
%% [NOTE: as printed, the first term of this display carries $x_2$ ("$a'_2x_2^{r_2+s_2}$") while every
%%  following term carries $x_3$; reproduced as printed.]
of the form (3.13).

If $\alpha + \frac{1}{x_3}\varphi(x_3)$ is not a unit series, then setting $\bar{y}_3 = \mu(y_3)y_3$, we obtain an expression (C) at $p_3$ with $s_m = 1$,
\begin{equation}
F = \tau z^m + a'_2 x_3^{r'_2} (\bar{y}_3 + \varphi'(x_3))^{s'_2} z^{m-2} + \cdots + a'_{m-1} x_3^{r'_{m-1}} (\bar{y}_3 + \varphi'(x_3))^{s'_{m-1}} z + \bar{a}_m x_3^{r'_m} \bar{y}_3).
\tag{c3}
\end{equation}
%% [NOTE: the display (c3) ends with a closing parenthesis that opens nowhere; reproduced as printed.]
As observed above, making the change of variables $\tilde{y}_3 = \bar{y}_3 + \varphi'(x_3)$ now gives a form (3.13).

If (b) holds at $p_3$, then an expression (2.1) holds at $p_3$ with a form (C) with $\operatorname{ord} \varphi < r$. If $\operatorname{ord} \varphi = 0$ we then have an expression (3.13) (with $s_i = 0$ for $2 \leqslant i \leqslant m-1$). If (c) holds then
\[
\bar{y}_2 + \varphi(x_2) = y_3 + \varphi(x_3 y_3) = y_3 \gamma(x_3, y_3)
\]
where $\gamma$ is a unit series. Thus we have an expression (2.2) with a form (3.13).

By induction on $r$, we must obtain an expression (2.1) or (2.2) with a form (3.13) at all points above $p$. We may assume that this already holds in $U_2$.

Suppose that $p_2$ is a 1-point above $p$.

%% --- page 8 (printed 2511) ---

Let
\[
J = I_2 \left( \begin{smallmatrix}
\frac{\partial u}{\partial x_2} & \frac{\partial u}{\partial y_2} & \frac{\partial u}{\partial z} \\[2pt]
\frac{\partial v}{\partial x_2} & \frac{\partial v}{\partial y_2} & \frac{\partial v}{\partial z}
\end{smallmatrix} \right) = x^n \left( \frac{\partial F}{\partial y_2}, \frac{\partial F}{\partial z} \right)
\]
for some positive integer $n$. Since $D$ contains the locus where $f$ is not smooth, we have that the localization $J_{\mathfrak{p}} = (\Oh{U_2,q})_{\mathfrak{p}}$, where $\mathfrak{p}$ is the prime ideal $(y_2, z)$ in $\Oh{U_2,q}$.
%% [NOTE: the localization is printed with the structure sheaf on the right-hand side
%%  ("$J_{\mathfrak p} = (\widehat{\mathcal O}_{U_2,q})_{\mathfrak p}$"); reproduced as printed.]

We compute
\[
\frac{\partial F}{\partial z} = \bar{a}_{m-1} x_2^{r_{m-1}} y_2^{s_{m-1}} + \Lambda_1 z
\]
and
\[
\frac{\partial F}{\partial y_2} = s_m \bar{a}_m x_2^{r_m} y_2^{s_m-1} + \Lambda_2 z
\]
for some $\Lambda_1, \Lambda_2 \in \Oh{U_2,q}$, to see that
\begin{equation}
\bar{a}_{m-1} \neq 0 \text{ and } s_{m-1} = 0, \text{ or } \bar{a}_m \neq 0 \text{ and } s_m = 1.
\tag{D}
\end{equation}

Suppose $p_2 \in U_2$ is a 2-point above $p$. Deforming $p_2$ to a 1-point, we see from (D) (and since (3.13) holds) that (3.10) holds at $p_2$.

Let $p_2 \in U_2$ be a 1-point above $p$ where the conclusions of the lemma do not hold. Let $T_3 \to T_2$ be the blow up of $\Lambda_2(p_2) \in T_2$ with induced blow up $\psi_3 : U_3 \to U_2$. Suppose $p_3 \in \psi_3^{-1}(p_2)$. We have that the conclusions of the lemma hold in the form (3.10) if $p_3$ is the 2-point which has permissible parameters $x_3, y_3, z$ defined by $x_2 = x_3 y_3$ and $y_2 = y_3$. At a 1-point which has permissible parameters $x_3, y_3, z$ defined by $x_2 = x_3$, $y_2 = x_3(y_3 + \alpha)$ with $\alpha \neq 0$, we have that a form (3.11) or (3.12) holds. Thus the only case where we may possibly have not achieved the conclusions of the lemma is at the 1-point which has permissible parameters $x_3, y_3, z$ defined by $x_2 = x_3$ and $y_2 = x_3 y_3$. We continue to blow up, so that there is at most one point $p_3$ above $p_2$ where the conclusions of the lemma do not hold. This point is a 1-point which has permissible parameters $x_3, y_3, z$ with $x_2 = x_3$ and $y_2 = x_3^n y_3$ where we can take $n$ as large as we like. Substituting into (3.13), we have an expression (2.1) at $p_3$ with
\[
F = \tau z^m + \bar{a}_2 x_3^{r_2+s_2 n} y_3^{s_2} z^{m-2} + \cdots + \bar{a}_{m-1} x_3^{r_{m-1}+s_{m-1} n} y_3^{s_{m-1}} z + \bar{a}_m x_3^{r_m+s_m n} y_3^{s_m}.
\]
Since we must have some $\bar{a}_i \neq 0$ with $s_i \leqslant 1$ by (D) for $n \gg 0$, we obtain a form (3.12) or (B). \hfill $\square$

Pages 882--883: The statement of Lemma 3.7 should be changed to:

%% =====================================================================
%% END OF THE COUNTED UNIT (statement + proof of Lemma 3.6).
%% =====================================================================

%% =====================================================================
%% AUXILIARY BLOCK --- NOT PART OF THE COUNTED UNIT.
%% The errata's own continuation of the same new case: the target form (F) of the revised
%% Lemma 3.7, and the inserted detailed proof that a point with a form (B) does have a form (F).
%% PDF pp. 10-12 = printed pp. 2513-2515. Transcribed for context only; the counted unit above
%% is complete without it.
%% =====================================================================

%% --- page 10 (printed 2513, top) ---

%% [The following is item (4) of the revised statement of Lemma 3.7, as printed at the top of the
%%  page; it is transcribed here only because the inserted proof below concludes with it.]
(4) $p_1^*$ is a 1-point, and we have an expression (2.1) with
\begin{equation}
F = \bar{\tau}_0 z^m + \bar{\tau}_2 \hat{x}_1^{r_2} \bar{y}_1 z^{m-2} + \cdots + \bar{\tau}_{m-1} \hat{x}_1^{r_{m-1}} \bar{y}_1 z + x_1^{r_m} y_1 + \hat{x}_1^t \bar{\Omega}
\tag{F}
\end{equation}
where $\bar{\tau}_0 \in \Oh{V_1,p_1^*}$ is a unit, $\bar{\tau}_i \in \Oh{V_1,p_1^*}$ are units (or zero) for $2 \leqslant i \leqslant m$, $\bar{\Omega} \in \Oh{V_1,p_1^*}$, $\bar{\tau}_i \neq 0$ for some $2 \leqslant i \leqslant m$ and $t > \omega'(m, r_2, \ldots, r_m)$. Further, $r_i > 0$ if $\bar{\tau}_i \neq 0$.

%% [Intervening errata item not transcribed: "Page 885, line 6: ..." (the correction of the
%%  hypothesis on $\bar\tau_m$).]

Page 885, line 11: after ``form (3.18)'' insert: ``and in the case when $p_1$ has a form (B) a similar argument shows that $p_1^*$ has a form (F).'' The following is a detailed proof of this statement.

Suppose that $p_1$ has a form (B). With the notation of Lemma 3.6, we have polynomials $\varphi, \psi$ such that
\[
x = \varphi(\tilde{x}_1, \tilde{y}_1), \quad y = \psi(\tilde{x}_1, \tilde{y}_1)
\]
determines the birational extension $\mathcal{O}_{T_0,p_0} \to \mathcal{O}_{T_1,\Lambda_1(p_1)}$, and we have a formal change of variables
\[
x_1 = \alpha(\tilde{x}_1, \tilde{y}_1)\tilde{x}_1, \quad y_1 = \beta(\tilde{x}_1, \tilde{y}_1)
\]
for some unit series $\alpha$ and series $\beta$. In this case, $p_1$ maps to a 1-point with an expression (3.13). Let $W \subset \{2, \ldots, m-1\}$ be the set of $i$ such that $s_i \geqslant 2$. Then $W \neq \emptyset$, $i \notin W$ implies, with the notation of (B),
\[
a_i(x, y) = \bar{a}_i(x_1, y_1) x_1^{r_i} y_1,
\]
\[
a_m(x, y) = \bar{a}_m(x_1, y_1) x_1^{r_m} y_1
\]
and
\[
\sum_{i \in W} a_i(x, y) z^i = x_1^t y_1^2 \Omega
\]
with $t > \omega'(m, r_1, \ldots, r_m)$.
%% [NOTE: the subscript is printed as $r_1$ here, where the rest of the errata writes
%%  $\omega'(m,r_2,\ldots,r_m)$; reproduced as printed.]

We have $x = \bar{\gamma} x^*$ with $\bar{\gamma} \equiv 1 \bmod m_p^r \Oh{X,p}$. Set $y^* = y$. At $p_1^*$, we have regular parameters $x_1^*, y_1^*$ in $\mathcal{O}_{T_1^*,\Lambda_1^*(p_1^*)}$ such that
\[
x^* = \varphi(x_1^*, y_1^*), \quad y^* = \psi(x_1^*, y_1^*),
\]
and $x_1^*, y_1^*, \tilde{z}$ are regular parameters in $\mathcal{O}_{V_1,\bar{p}_1^*}$ (recall that $z = \sigma \tilde{z}$ in Lemma 3.1). We have regular parameters $\bar{x}_1, \bar{y}_1, \in \Oh{T_1^*,\Lambda_1^*(p_1^*)}$ defined by
\[
\bar{x}_1 = \alpha(x_1^*, y_1^*) x_1^*, \quad \bar{y}_1 = \beta(x_1^*, y_1^*).
\]
%% [NOTE: printed with a comma after $\bar y_1$ and before the membership sign; reproduced as printed.]

We have $u = x^a = x_1^{a_1}$ where $a_1 = ad$ for some $d \in \mathbb{Z}_+$. Since $[\alpha(\tilde{x}_1, \tilde{y}_1)\tilde{x}_1]^d = x$, we have that $[\alpha(x_1^*, y_1^*)x_1^*]^d = x^*$. Set $\hat{x}_1 = \bar{\gamma}^{\frac{1}{d}} \bar{x}_1 = \bar{\gamma}^{\frac{1}{d}} \alpha(x_1^*, y_1^*) x_1^*$. We have that $\bar{\gamma}^{\frac{1}{d}} \alpha(x_1^*, y_1^*)$ is a unit in $\Oh{V_1,p_1^*}$, and $x = \hat{x}_1^d$.

%% --- page 11 (printed 2514) ---

Thus $x_1 = \hat{x}_1$ (with an appropriate choice of root $\bar{\gamma}^{\frac{1}{d}}$). We have $u = \hat{x}_1^{ad}$, so that $\hat{x}_1, \bar{y}_1, z$ are permissible parameters at $p_1^*$.

For $i \notin W$, we have
\[
a_i(x, y) = a_i(\bar{\gamma} x^*, y^*) \equiv a_i(x^*, y^*) \bmod m_p^r \Oh{V,p}
\]
and
\begin{align*}
a_i(x^*, y^*) &= a_i(\varphi(x_1^*, y_1^*), \psi(x_1^*, y_1^*)) \\
&= \bar{x}_1^{r_i} \bar{a}_i(\bar{x}_1, \bar{y}_1) \\
&\equiv x_1^{r_i} \bar{a}_i(x_1, \bar{y}_1) \bmod m_p^r \mathcal{O}_{V_1,p_1^*}.
\end{align*}
We further have
\[
a_m(x^*, y^*) \equiv x_1^{r_m} \bar{y}_1 \bmod m_p^r \Oh{V_1,p_1^*}
\]
and
\[
x_1^t y_1^2 \Omega \equiv x_1^t \bar{y}_1^2 \Omega(x_1, \bar{y}_1, z) \bmod m_p^r \mathcal{O}_{V_1,p_1^*}.
\]
Thus we have expressions (2.1) with
\[
u = x_1^{da}
\]
\begin{equation}
v = P(x_1^d) + x_1^{bd} P_1(x_1) + x_1^{bd} (\bar{\tau} z^m + x_1^{r_2} \bar{a}_2(x_1, \bar{y}_1)\bar{y}_1 z^{m-2} + \cdots + x_1^{r_m} \bar{y}_1 + x_1^t \bar{y}_1^2 \bar{\Omega} + h)
\tag{c4}
\end{equation}
where $\bar{\tau} \in \Oh{V_1,p_1^*}$ is a unit series and
\[
h \in m_p^r \Oh{V_1,p_1^*} \subset (x_1, z)^r.
\]
Set $s = r - m$, and write
\[
h = z^m \Lambda_0(x_1, \bar{y}_1, z) + z^{m-1} x_1^{1+s} \Lambda_1(x_1, \bar{y}_1) + z^{m-2} x_1^{2+s} \Lambda_2(x_1, \bar{y}_1) + \cdots + z x_1^{(m-1)+s} \Lambda_{m-1}(x_1, \bar{y}_1) + x_1^{m+s} \Lambda_m(x_1, \bar{y}_1)
\]
with $\Lambda_0 \in m_{p_1^*} \Oh{V_1,p_1^*}$ and $\Lambda_i \in \Bbbk[[x_1, \bar{y}_1]]$ for $1 \leqslant i \leqslant m$.

Substituting into (c4), we obtain an expression
\[
u = x_1^{da}
\]
\[
v = P(x_1^d) + x_1^{bd} P_1(x_1) + x_1^{bd} (\bar{\tau}_0 z^m + x_1^{r_2} \bar{\tau}_2 \bar{y}_1 z^{m-2} + \cdots + x_1^{r_{m-1}} \bar{\tau}_{m-1} \bar{y}_1 z + x_1^{r_m} \bar{y}_1 + x_1^{t'} \bar{\Omega})
\]
where $\bar{\tau}_0 \in \Oh{V_1,p_1^*}$ is a unit, $\bar{\tau}_i \in \Oh{V_1,p_1^*}$ are units (or zero), for $1 \leqslant i \leqslant m-1$.

%% --- page 12 (printed 2515) ---

We have $\bar{\tau}_0 = \tau + \Lambda_0$, $\tau_i = \bar{a}_i(x_1, \bar{y}_1)$ for $2 \leqslant i \leqslant m-1$, and
\[
\bar{\Omega} = x_1^{t-t'} \bar{y}_1^2 \Omega(x, \bar{y}_1, z) + z^{m-1} x_1^{1+s-r_m} \Lambda_1(x_1, \bar{y}_1) + \cdots + x_1^{m+s-r_m} \Lambda_m(x_1, \bar{y}_1)).
\]
%% [NOTE: printed (i) with $\tau_i$ unbarred on the left of the first line, and (ii) with a doubled
%%  closing parenthesis at the end of the display of $\bar\Omega$; reproduced as printed.]
with $t' = \omega'(m, r_2, \ldots, r_m) + 1$

We thus have the desired form (F).

%% =====================================================================
%% END OF AUXILIARY BLOCK.
%% =====================================================================

\end{document}
